%======================================================================
% Section 12: Comparative Synthesis of Existing Research
%======================================================================
\section{Comparative Synthesis of Existing Research}
\label{sec:comparative}

This section synthesizes the existing research across all theory families using a twelve-criteria comparison framework. The synthesis is constructive: differences between theories are presented as modeling choices with different assumptions and domains of validity, not as contradictions or errors.

\subsection{Cross-theory comparison table (12 criteria $\times$ 9 theory families)}
\label{sec:12.1}

\begin{table}[h]
\centering
\small
\caption{Twelve-criteria comparison across theory families. Flux status: V=verified, P=partial, U=uncertain.}
\label{tab:twelve_criteria}
\begin{tabularx}{\textwidth}{@{}lcccccccccc@{}}
\toprule
\textbf{Criterion} & \textbf{CD} & \textbf{LS} & \textbf{GL} & \textbf{GN-II} & \textbf{GN-III} & \textbf{DPL} & \textbf{TPL} & \textbf{MGT} & \textbf{MDD} & \textbf{Nonlocal} \\
\midrule
Relaxation parameters & 0 & 1 & 2 & 0 & 1 & 2 & 3 & 1 & 1+kernel & $(e_0a)^2$ \\
Dissipation & Yes & Yes & No & No & Yes & Yes & Yes & Yes & Kernel-dep. & Yes \\
Wave speed & Infinite & Finite & Finite & Finite & Finite & Finite & Finite & Finite & Finite & Finite (dispersive) \\
Energy flux & V & P & U & U & U & P & U & U & U & U \\
Well-posedness & Yes & $\tau_0>0$ & Yes & Yes & Yes & $\tau_q\geq\tau_T$ & Yes & Yes & Kernel-dep. & Yes \\
Causality & No & Yes & Yes & Yes & Yes & Yes & Yes & Yes & Yes & Yes \\
Interface conditions & Classical & Classical & Classical & Classical & Classical & Classical & Classical & Classical & Classical & Modified \\
Mathematical structure & PDE & PDE & PDE & PDE & PDE & PDE & PDE & 3rd-order PDE & Integro-PDE & Integro-PDE \\
Primary source & Biot \cite{Biot1956} & LS \cite{Lord1967} & GL \cite{Green1972} & GN \cite{Green1991} & GN \cite{Green1993} & Tzou \cite{Tzou1995} & Tzou \cite{Tzou1997} & Quintanilla \cite{Quintanilla2019} & Wang \& Li 2011 & Eringen \cite{Eringen1972} \\
Open access & No & No & Yes & No & No & No & No & No & Yes (PMC) & No \\
Applications & General & Thermal shocks & Symmetric bodies & Conservative & General & Micro-scale & Micro-scale & Thermal waves & Memory materials & Nanoscale \\
\bottomrule
\end{tabularx}
\end{table}

\subsection{Evidence trends by theory family}
\label{sec:12.2}

\begin{itemize}
\item \textbf{Classical (CD):} Most extensively validated experimentally; infinite speed paradox well documented.
\item \textbf{LS:} Widely applied to thermal shock problems; relaxation time parameter well understood.
\item \textbf{GL:} Applied to centrally symmetric bodies; dissipationless in appropriate limits.
\item \textbf{GN-II:} Conservative theory; no dissipation; thermal waves without attenuation.
\item \textbf{GN-III:} Most general GN model; combines dissipative and non-dissipative characteristics.
\item \textbf{DPL:} Applied to micro-scale heat transfer; two phase lags provide modeling flexibility.
\item \textbf{TPL:} Extension of DPL with thermal displacement gradient; less widely applied.
\item \textbf{MGT:} Third-order equation; thermal inertia; emerging theory.
\item \textbf{MDD:} Integral-form memory; alternative to fractional calculus; growing literature.
\item \textbf{Nonlocal:} Applied to nanoscale materials; size-dependent effects well demonstrated.
\end{itemize}

\begin{figure}[htbp]
\centering
\includegraphics[width=0.9\textwidth]{fig07_theory_relationships.pdf}
\caption{Relationships and limiting cases between thermoelastic theory families. All generalized theories converge to classical coupled thermoelasticity as relaxation parameters approach zero.}
\label{fig:theory_relationships_sec12}
\end{figure}

\subsection{Mathematical relationships: limiting cases and specializations}
\label{sec:12.3}

Table \ref{tab:limiting_cases_revisited} summarizes the limiting relationships between theories.

\begin{table}[h]
\centering
\small
\caption{Limiting cases between thermoelastic theories.}
\label{tab:limiting_cases_revisited}
\begin{tabularx}{\textwidth}{@{}llX@{}}
\toprule
\textbf{From} & \textbf{To} & \textbf{Conditions} \\
\midrule
LS & Classical & $\tau_0 \to 0$ \\
DPL & Classical & $\tau_q = \tau_T = 0$ \\
TPL & DPL & $\tau_\nu = 0$ \\
MDD+TPL & Classical & All $\tau \to 0$, kernel $\to 1$ \\
MDD+TPL & LS & $\tau_T = \tau_\nu = 0$, first order \\
MDD+TPL & DPL & $\tau_\nu = 0$ \\
GN-II & Classical & Special parameter limit \\
\bottomrule
\end{tabularx}
\end{table}

\subsubsection{Note on limiting procedures}

All limiting procedures in the existing literature are parameter substitutions, not rigorous limit theorems with convergence proofs or error bounds. This is documented as a limitation of the current mathematical framework, not as an error in any publication.

\subsection{Contradiction mapping: where studies disagree and why}
\label{sec:12.4}

Apparent contradictions between studies arise from:

\begin{enumerate}
\item \textbf{Different assumptions:} Constitutive assumptions, material symmetry, geometry, boundary conditions differ between studies.
\item \textbf{Different approximation levels:} First-order vs. higher-order Taylor expansions of phase-lag laws.
\item \textbf{Different normalization:} Different choices of reference temperature, density, or elastic constants.
\item \textbf{Different frequency conventions:} $e^{i\omega t}$ vs. $e^{-i\omega t}$ affect sign conventions.
\end{enumerate}

Each apparent contradiction is traceable to specific assumptions, not to mathematical inconsistency. No study is presented as erroneous; differences reflect different modeling choices.

\subsection{Model comparability: conditions under which theories can be compared}
\label{sec:12.5}

Two theories can be meaningfully compared when:

\begin{enumerate}
\item They share the same geometric setting (e.g., both half-spaces, both plates),
\item They use the same boundary and interface conditions,
\item They use the same normalization and units,
\item They address the same physical phenomenon (e.g., wave reflection at an interface),
\item The comparison accounts for differences in relaxation parameters and kernel functions.
\end{enumerate}

\subsection{Notation summary for Section 12}

\begin{table}[h]
\centering
\small
\caption{Notation introduced in Section 12.}
\label{tab:notation-sec12}
\begin{tabularx}{\textwidth}{@{}lXlr@{}}
\toprule
\textbf{Symbol} & \textbf{Meaning} & \textbf{Units} & \textbf{Equation} \\
\midrule
CD & Classical coupled theory & --- & Table \ref{tab:twelve_criteria} \\
LS & Lord--Shulman theory & --- & Table \ref{tab:twelve_criteria} \\
GL & Green--Lindsay theory & --- & Table \ref{tab:twelve_criteria} \\
GN-II/III & Green--Naghdi type II/III & --- & Table \ref{tab:twelve_criteria} \\
DPL/TPL & Dual/triple phase-lag & --- & Table \ref{tab:twelve_criteria} \\
MGT & Moore--Gibson--Thompson & --- & Table \ref{tab:twelve_criteria} \\
MDD & Memory-dependent derivative & --- & Table \ref{tab:twelve_criteria} \\
\bottomrule
\end{tabularx}
\end{table}

\subsection{Open questions for Section 12}

\begin{enumerate}[label=\textbf{OQ12.\arabic*}]
\item Can rigorous limit theorems be established between all theory pairs? \label{oq:12.1}
\item What is the minimal set of criteria needed to distinguish theory families? \label{oq:12.2}
\item How should apparent contradictions between studies be systematically classified? \label{oq:12.3}
\end{enumerate}

% LocalWords: 
